% FIGURE NEEDED: flaring disk cross-section with the star, surface layer, tangent line, angles h/r, dh/dr, beta and height h(r) at radius r (2022 Solutions, p. 3)
\begin{description}
\item[(a)] The angle between the light beam and the horizontal plane is
  \[
    \theta_l \approx \tan\theta_l = \frac{h(r)}{r}
  \]
  \hfill(1.0 points)

  Additionally, a tangent is inclined to the disc plane with an angle of
  \[
    \theta_T \approx \tan\theta_T = \frac{\Delta h(r)}{\Delta r}
  \]
  \hfill(2.0 points)

  An exterior angle of a triangle is equal to the sum of its two interior
  opposite angles. Thus,
  \[
    \beta = \theta_T - \theta_l = \frac{\Delta h(r)}{\Delta r} - \frac{h(r)}{r}
  \]
  \hfill(1.0 points)
\item[(b)] The flux from of stellar radiation at distance $r$ from the star is
  $E_s = \frac{L_s}{4\pi r^2}$, where $L_s$ is the luminosity of the star.
  However, the irradiation flux is the normal projection of this flux onto the
  infinitesimal small surface $A$ of the disk.
  \[
    Q_+ = E_s A\sin\beta \approx \frac{L_s A\beta}{4\pi r^2}
  \]
  \hfill(1.5 points)

  From the Stefan-Boltzmann law, the cooling rate is give by
  \[
    Q_- = A\sigma T_D^4.
  \]
  \hfill(0.5 points)

  In the thermal equilibrium
  \[
    Q_+ = Q_- \Longrightarrow
    T_D = \left(\frac{L_s\beta}{4\pi\sigma r^2}\right)^{\!\frac{1}{4}}
  \]
  \hfill(1.0 points)
\item[(c)] The condition for the isothermal layer can be solved by applying
  \begin{align*}
    \beta &= \frac{\Delta h(r)}{\Delta r} - \frac{h(r)}{r}
      = \frac{a(r+\Delta r)^b - ar^b}{\Delta r} - \frac{ar^b}{r} \\
    &= ar^b\left(\frac{(1+\frac{\Delta r}{r})^b - 1}{\Delta r} - \frac{1}{r}\right)
      = ar^b\left(\frac{b\Delta r}{r\Delta r} - \frac{1}{r}\right) \\
    &= a(b-1)r^{b-1} \propto r^2 \\
    &\Longrightarrow b = 3
  \end{align*}
  \hfill(2.0 points)

  Finally, one can determine the second constant by
  \begin{align*}
    \beta &= a(b-1)r^{b-1} = 2ar^2 \\
    \therefore\; T_{SL}^4 &= \left(\frac{aL_s}{2\pi\sigma}\right) \\
    a &= \left(\frac{2\pi\sigma T_{SL}^4}{L_s}\right)
  \end{align*}
  \hfill(1.0 points)
\end{description}
